\begin{abstract}
A language model that learns tasks through adapters forgets, because a new adapter changes what the
model computes on old inputs. FK-LoRA gives every task one frozen set of adapters and sends each input
to one of them, so an answer changes only if its route changes. Our main rule,
\recog{}, reads no instruction: it is a recogniser on frozen model states, does not depend on the order of arrival and blocks no task. It comes with an audit, which estimates and bounds from held-out inputs the
share of changed answers, and a variant, \handicap{}, which bounds in advance what a later task
can take. Without instructions any method loses on two tasks together at least the share of inputs that they have in common with different
right answers. Benchmarks hide this, as their inputs begin with the instruction. Without it, the two
published methods we ran, GainLoRA and InfLoRA, change \AudSeedsGainloraChanged{} and
\AudSeedsInfloraChanged{} of \AudSeedsCompared{} answers (\FiveGainChanged{} and \FiveInfChanged{}
with it). This loss was not forced: no input of these benchmarks occurs in two tasks, and the recogniser gives \AudDevLLamOneEnd{}\% of test inputs to their task on 15 development tasks,
\AudScaleLongThreeEnd{}\% on a second benchmark and \AudFreshFourthLEnd{}\% on 30 unseen tasks. On five tasks (three seeds, a tenth of the published epochs) plain adapters behind it change \AudSeedsRidgeChanged{} answers and end at
\AudSeedsRidgeAPOne{} Rouge-L, against \AudSeedsGainloraAPOne{} and \AudSeedsInfloraAPOne{}; most of that gap is learning. With the published 100 epochs (one seed) the two methods learn the tasks as well as plain adapters and
then forget \AudFullGainloraFTInt{} and \AudFullInfloraFTInt{} points, ending at
\AudFullGainloraAPOne{} and \AudFullInfloraAPOne{}, against \AudFullOursRidgeAPOne{} for plain
adapters behind it (\AudOursRidgeAPOne{} at 10 epochs); with their instructions they forget \AudFullDefOriginalGainloraFTOne{} and \AudFullDefOriginalInfloraFTOne{} points. When inputs
carry an instruction, a second rule gives each task a zone around the word-vectors of its
instruction, and we prove that an answer inside a zone never changes; no answer to an earlier task changed (\FiveCertChanged{} of \FiveCertCompared{}
on five tasks, three seeds; \FifteenCertChanged{} of \FifteenCertCompared{} on fifteen). Reworded instructions leave the zones: on machine paraphrases of
\WidenTasksEval{} real tasks they certify \WidenRefTfivePLow{}\% to \WidenRefTfivePHigh{}\% of the
inputs for the right task, and five attempts to widen them missed their bars. The recogniser ignores the instruction, so no rewording changes its routes: with instructions (five tasks, three seeds, development rewordings) it ends at \AudUThreeHeavyRidgeAPOne{} under a heavy rewording where the zones end at \AudUThreeHeavyCertAPOne{}, but with the original instructions the zones end higher (\AudUThreeOriginalCertAPOne{} against \AudUThreeOriginalRidgeAPOne{}) and change no answer where the recogniser changes \AudUThreeOriginalRidgeChanged{} of \AudUThreeCompared{}. Without an instruction, inputs that occur in two tasks can at best be split between them; the bound in advance is loose; all runs use T5-base, and nothing is in the published setting. Stand-alone code is released.
\end{abstract}
